% !PRINTED 393
\weberpage{393}
%<<CONT
2 haben, noch wegzulassen sind, ergeben sich die reducirten
Zahlen:
\[
\{5, 4, 5\}, \quad \{4, 6, 5\}, \quad \{5, 6, 4\}, \quad \{1, 10, 4\}, \quad \{4, 10, 1\},
\]
die sich so in eine Periode ordnen, wobei das erste Glied am
Ende noch einmal zugesetzt ist:
\[
\{4, 10, 1\}, \quad \{1, 10, 4\}, \quad \{4, 6, 5\}, \quad \{5, 4, 5\}, \quad \{5, 6, 4\}, \quad \{4, 10, 1\},
\]
und für $\omega$ ergiebt sich die Kettenbruchperiode:
\[
5 + \sqrt{29} = [10, 2, 1, 1, 2].
\]
Wenn man mit dem vierten Gliede anfängt, also den Kettenbruch
$[1, 2, 10, 2, 1]$ betrachtet, der die Entwickelung von $\dfrac{2 + \sqrt{29}}{5}$
ist, so sieht man, dass die Periode \gesperrt{umkehrbar} ist.

3. $D = 76 = 4 \ldotp 19, \quad \lambda = 8$
% [resolved 2026-08-06, 5x zoom strip 2/5] Confirmed. The line reads "3. D = 76 = 4. 19,
% lambda = 8": the mark is Weber's raised multiplication period, and the arithmetic settles
% it -- 4 x 19 = 76. Transcription stands. Same strip confirms b = 2, 4, 6, 8 and
% ac = 18, 15, 10, 3.
\[
b = \quad 2, \quad 4, \quad 6, \quad 8
\]
\[
a\,c = 18, \quad 15, \quad 10, \quad 3.
\]
\gesperrt{Grenzen} für $a$ und $c$:
\[
4, 5; \quad 3, 6; \quad 2, 7; \quad 1, 8;
\]
also sind die reducirten Zahlen:
\[
\{3, 4, 5\}, \quad \{5, 4, 3\}, \quad \{2, 6, 5\}, \quad \{5, 6, 2\}, \quad \{1, 8, 3\}, \quad \{3, 8, 1\}.
\]
Wir erhalten eine Periode:
\[
\{3, 8, 1\}, \quad \{1, 8, 3\}, \quad \{3, 4, 5\}, \quad \{5, 6, 2\}, \quad \{2, 6, 5\}, \quad \{5, 4, 3\}, \quad \{3, 8, 1\}
\]
und die Periode des Kettenbruches für $4 + \sqrt{19}$ wird
\[
[8, 2, 1, 3, 1, 2].
\]
Auch dieser Kettenbruch hat eine umkehrbare Periode, da
man dieselbe Reihenfolge der Zahlen erhält, mag man von der
ersten 2 nach rechts oder von der zweiten 2 nach links lesen.

4. $D = 37, \qquad \lambda = 6,$
\[
b = 1, \quad 3, \quad 5
\]
\[
a\,c = 9, \quad 7, \quad 3;
\]
\gesperrt{Grenzen} für $a$, $c$:
\[
3, 3; \quad 2, 4; \quad 1, 5;
\]
\gesperrt{reducirte Zahlen}:
\[
\{3, 1, 3\}, \quad \{1, 5, 3\}, \quad \{3, 5, 1\};
\]
\gesperrt{Periode}:
\[
\{3, 5, 1\}, \quad \{1, 5, 3\}, \quad \{3, 1, 3\}, \quad \{3, 5, 1\};
\]
\gesperrt{Kettenbruch} mit umkehrbarer Periode $[1, 5, 1]$:
%CONT>>
